% =====================================================================
% Section 6: Conclusion — EvidenceFlow
% =====================================================================
%
% Status: Draft Section 6 Conclusion v1 (2026-05-17)
% Length: ~108 words, single paragraph, 5 sentences. No bullet list,
%         no headline numbers, no new terminology.
% Style reference: composite of NeocorRAG (problem diagnosis) +
%                  Relink (principle close) + HGRAG (compact module
%                  recap), per advisor feedback on conclusion length
%                  and registers in paper/refer/.
%
% =====================================================================

\section{Conclusion}
\label{sec:conclusion}

We presented \methodname{}, a graph-based retrieval framework that
treats graph propagation as the start of evidence selection rather
than the final ranking. \methodname{} builds a fact-source substrate
that promotes each OpenIE fact to a graph node and binds it to its
source passage, and runs query-local personalized PageRank on an
induced fact-and-passage graph to produce a query-conditioned
evidence flow. Two evidence-aware steps then refine this flow:
relation-grounded edges to admit bridge passages whose relations
phrases match the question, and coverage-aware reranking selects a
list that covers distinct question-side requirements. Experiments
on three multi-hop QA benchmarks show consistent gains over dense,
OpenIE-based, and chain-oriented baselines, with ablations
isolating the contribution of each operator and a substrate-isolation
experiment showing that fact-source grounding mainly contributes to
connected evidence coverage on bridge-heavy questions. The results suggest
that converting propagated relevance into a connected and covering
evidence list is a productive design axis for GraphRAG.
